\section{Node-Death Detection and Reassignment}\label{sec:death}
No assigned chunk may be silently lost, or the search could terminate reporting
``not found'' while skipping the key. The master defends against two failure
modes, both of which requeue the dead node's chunk to the \emph{front} of
\texttt{pending} so a surviving worker picks it up:

\begin{enumerate}\itemsep2pt
\item \textbf{Hard crash / kill.} The TCP socket closes; the master's read
returns EOF, and the disconnect handler requeues every chunk that worker held.
\item \textbf{Freeze / network partition.} The socket stays open but
\texttt{PROGRESS} heartbeats stop. A monitor thread requeues any chunk whose
last heartbeat is older than a timeout $T$ (default $5$\,s) and marks the node
dead.
\end{enumerate}

\paragraph{Empirical check.} In a controlled run with three workers, freezing
one with \texttt{SIGSTOP} and killing another with \texttt{SIGKILL} produced the
expected master events:
\begin{quote}\footnotesize\ttfamily
TIMEOUT: worker freeze-w1 silent >4.0s, requeued chunk 19\\
worker kill-w2 disconnected; requeued chunks [40]
\end{quote}
Both chunks were subsequently completed by the surviving worker, so coverage of
the keyspace remained total.

% TODO(you): optional reflection on timeout tuning (false positives on slow
% nodes vs. detection latency). Left for your prose.
\todo{optional: discuss timeout tuning trade-offs}
